% related_reasonspec.tex — Related Work for the Reasoning Speculation paper
% Loaded by main_reasonspec.tex via % related_reasonspec.tex — Related Work for the Reasoning Speculation paper
% Loaded by main_reasonspec.tex via % related_reasonspec.tex — Related Work for the Reasoning Speculation paper
% Loaded by main_reasonspec.tex via % related_reasonspec.tex — Related Work for the Reasoning Speculation paper
% Loaded by main_reasonspec.tex via \input{sections/related_reasonspec}

% ---------------------------------------------------------------------------
\subsection{Test-Time Compute Scaling}
\label{sec:related:ttc}
% ---------------------------------------------------------------------------

The observation that ``thinking longer'' improves accuracy has motivated a new class of reasoning-specialized LLMs.
OpenAI o1~\citep{openai2024o1} first demonstrated that training models for extended chain-of-thought generation yields substantial gains on mathematical and scientific benchmarks.
QwQ~\citep{qwq2024} and DeepSeek-R1~\citep{deepseekr1} further showed that reinforcement-learning-based training for long reasoning produces strong open-weight reasoners.
\citet{snell2024scaling} formalized this as a \emph{test-time compute scaling law}: accuracy improves log-linearly with token budget on mathematical tasks.
However, this scaling is far from free---\citet{chen2024overthinking} and our own analysis (\S\ref{sec:intro}) show that accuracy \emph{degrades} on a non-trivial fraction of instances when the budget is increased, a phenomenon termed ``overthinking.''
Budget forcing~\citep{muennighoff2025s1} represents an early attempt at controlling reasoning length by truncating the thinking trace at a fixed token limit and appending a wrap-up prompt.
While effective as a diagnostic tool, budget forcing applies a uniform budget to all instances and thus cannot exploit the heterogeneity of problem difficulty---precisely the gap \method{} is designed to fill.
The ``sweet spot'' analysis of~\citet{chen2024overthinking} confirms that every instance has an optimal budget, but no fixed allocation can simultaneously serve easy and hard problems.

% ---------------------------------------------------------------------------
\subsection{Adaptive Budget Allocation}
\label{sec:related:adaptive}
% ---------------------------------------------------------------------------

A growing literature seeks to allocate reasoning compute \emph{per instance}, which we organize into three families.

\paragraph{Training-based methods.}
Several approaches modify the model or its training objective to produce difficulty-calibrated traces.
AdaThink~\citep{adathink2025} shapes rewards via GRPO to encourage the model to terminate reasoning early when confident, achieving budget savings without retraining from scratch.
AdaToken~\citep{adatoken2025} frames budget allocation as an online bandit problem and provides regret bounds, but requires online interaction during inference.
SelfBudgeter~\citep{liu2025selfbudgeter} trains the model to predict its own required budget before generation, enforcing adherence via budget-guided GRPO.
TALE~\citep{han2025tale} includes the target token budget in the prompt and trains via SFT/DPO to compress reasoning.
While these methods can be effective, they require model fine-tuning and are thus coupled to specific model weights and training data---a significant deployment constraint that \method{} avoids entirely.

\paragraph{Confidence-based methods.}
An alternative family uses the model's self-reported or estimated confidence to trigger early exit.
Satisficing Reasoning~\citep{luo2025satisficing} applies a confidence threshold to halt generation when the model deems its current answer sufficient.
Certaindex~\citep{luo2025certaindex} and Dynasor~\citep{fu2025dynasor} monitor uncertainty during generation and route to shorter or longer reasoning accordingly.
These methods are training-free but rely on a \emph{single-path} confidence signal, which we show is fundamentally limited: the model's self-assessed confidence is poorly calibrated, and single-trace uncertainty estimates achieve Spearman $|\rho| < 0.2$ with true problem difficulty (\S\ref{sec:theory}).

\paragraph{Training-free injection methods.}
RTB~\citep{hou2025thinkless} and Think Less Think Better~\citep{hou2025thinkless} inject wrap-up tokens (e.g., ``\texttt{Wait, let me reconsider\ldots}'') into the reasoning trace to steer the model toward shorter or longer deliberation without any training.
In our comprehensive evaluation (\S\ref{sec:intro}), we tested all major categories of single-path controllers---including feature-based, uncertainty-based, and injection-based methods---and found that \textbf{none outperformed the fixed-budget baseline}, with degradations of 5--7\,pp.
The root cause is informational: \emph{any} difficulty signal extracted from a single reasoning trace is bounded in mutual information with problem difficulty by $I(D; X_1)$ (Proposition~\ref{prop:single-path}), which is strictly less than the consensus signal $I(D; C_K)$ for $K \ge 2$.
This negative result motivates \method{}'s fundamental departure: rather than extracting a weak signal from one path, restructure the computation to produce a multi-path signal.

% ---------------------------------------------------------------------------
\subsection{Self-Consistency and Multi-Path Reasoning}
\label{sec:related:sc}
% ---------------------------------------------------------------------------

Self-consistency~\citep{wang2023selfconsistency} improves accuracy by generating $K$ complete reasoning traces and returning the majority-vote answer.
Best-of-$N$ with reward models~\citep{cobbe2021gsm8k,lightman2023prm} samples $N$ solutions and selects the highest-scoring one according to a learned verifier.
MCTS-based reasoning methods~\citep{hao2023reasoning,xie2024selfevaluation,qi2024mutual} build search trees over reasoning steps, using rollout-based value estimates to guide expansion.
These approaches all use multiple paths to improve \emph{answer quality}---but they allocate compute \emph{uniformly} across instances, spending the same $K$ full-budget traces on every question regardless of difficulty.
\method{} draws on the same insight that multi-path agreement correlates with correctness, but uses consensus as a \emph{difficulty signal for adaptive routing}, not merely as an aggregation mechanism.
Concretely, \method{} runs $K$ \emph{short} probes (not full-budget traces), uses the agreement pattern to \emph{decide} whether additional compute is warranted, and invests that compute \emph{selectively} on hard instances.
This yields a fundamentally different cost profile: self-consistency pays $K \times B_{\mathrm{full}}$ tokens per question; \method{} pays $K \times B_{\mathrm{probe}}$ for easy questions and scales up only where consensus indicates difficulty.

% ---------------------------------------------------------------------------
\subsection{Speculative Decoding}
\label{sec:related:spec}
% ---------------------------------------------------------------------------

Speculative decoding~\citep{leviathan2023fast,chen2023accelerating} accelerates autoregressive generation by using a small, fast draft model to propose token sequences that a larger target model then verifies in parallel.
The key insight is that \emph{restructuring the computation}---from sequential single-model generation to a draft-then-verify pipeline---can reduce latency without changing the output distribution.
This paradigm has been extended to tree-structured speculation~\citep{miao2024specinfer}, self-speculation within a single model~\citep{zhang2024draft}, and cascade routing across model sizes~\citep{chen2023frugalgpt}.

\method{} draws a deliberate analogy to speculative decoding, applied at the level of \emph{reasoning} rather than token generation.
Where speculative decoding uses a cheap draft model to identify tokens the target model would likely accept, \method{} uses cheap parallel probes to identify questions the model can likely answer without extended deliberation.
Where speculative decoding falls back to the target model for tokens the draft model gets wrong, \method{} falls back to full deliberation for questions where the probes disagree.
In both cases, the approach exploits a structural asymmetry---most tokens are easy to predict, most questions are easy to solve---to reduce expected cost while preserving or improving quality, without modifying the underlying model.


% ---------------------------------------------------------------------------
\subsection{Test-Time Compute Scaling}
\label{sec:related:ttc}
% ---------------------------------------------------------------------------

The observation that ``thinking longer'' improves accuracy has motivated a new class of reasoning-specialized LLMs.
OpenAI o1~\citep{openai2024o1} first demonstrated that training models for extended chain-of-thought generation yields substantial gains on mathematical and scientific benchmarks.
QwQ~\citep{qwq2024} and DeepSeek-R1~\citep{deepseekr1} further showed that reinforcement-learning-based training for long reasoning produces strong open-weight reasoners.
\citet{snell2024scaling} formalized this as a \emph{test-time compute scaling law}: accuracy improves log-linearly with token budget on mathematical tasks.
However, this scaling is far from free---\citet{chen2024overthinking} and our own analysis (\S\ref{sec:intro}) show that accuracy \emph{degrades} on a non-trivial fraction of instances when the budget is increased, a phenomenon termed ``overthinking.''
Budget forcing~\citep{muennighoff2025s1} represents an early attempt at controlling reasoning length by truncating the thinking trace at a fixed token limit and appending a wrap-up prompt.
While effective as a diagnostic tool, budget forcing applies a uniform budget to all instances and thus cannot exploit the heterogeneity of problem difficulty---precisely the gap \method{} is designed to fill.
The ``sweet spot'' analysis of~\citet{chen2024overthinking} confirms that every instance has an optimal budget, but no fixed allocation can simultaneously serve easy and hard problems.

% ---------------------------------------------------------------------------
\subsection{Adaptive Budget Allocation}
\label{sec:related:adaptive}
% ---------------------------------------------------------------------------

A growing literature seeks to allocate reasoning compute \emph{per instance}, which we organize into three families.

\paragraph{Training-based methods.}
Several approaches modify the model or its training objective to produce difficulty-calibrated traces.
AdaThink~\citep{adathink2025} shapes rewards via GRPO to encourage the model to terminate reasoning early when confident, achieving budget savings without retraining from scratch.
AdaToken~\citep{adatoken2025} frames budget allocation as an online bandit problem and provides regret bounds, but requires online interaction during inference.
SelfBudgeter~\citep{liu2025selfbudgeter} trains the model to predict its own required budget before generation, enforcing adherence via budget-guided GRPO.
TALE~\citep{han2025tale} includes the target token budget in the prompt and trains via SFT/DPO to compress reasoning.
While these methods can be effective, they require model fine-tuning and are thus coupled to specific model weights and training data---a significant deployment constraint that \method{} avoids entirely.

\paragraph{Confidence-based methods.}
An alternative family uses the model's self-reported or estimated confidence to trigger early exit.
Satisficing Reasoning~\citep{luo2025satisficing} applies a confidence threshold to halt generation when the model deems its current answer sufficient.
Certaindex~\citep{luo2025certaindex} and Dynasor~\citep{fu2025dynasor} monitor uncertainty during generation and route to shorter or longer reasoning accordingly.
These methods are training-free but rely on a \emph{single-path} confidence signal, which we show is fundamentally limited: the model's self-assessed confidence is poorly calibrated, and single-trace uncertainty estimates achieve Spearman $|\rho| < 0.2$ with true problem difficulty (\S\ref{sec:theory}).

\paragraph{Training-free injection methods.}
RTB~\citep{hou2025thinkless} and Think Less Think Better~\citep{hou2025thinkless} inject wrap-up tokens (e.g., ``\texttt{Wait, let me reconsider\ldots}'') into the reasoning trace to steer the model toward shorter or longer deliberation without any training.
In our comprehensive evaluation (\S\ref{sec:intro}), we tested all major categories of single-path controllers---including feature-based, uncertainty-based, and injection-based methods---and found that \textbf{none outperformed the fixed-budget baseline}, with degradations of 5--7\,pp.
The root cause is informational: \emph{any} difficulty signal extracted from a single reasoning trace is bounded in mutual information with problem difficulty by $I(D; X_1)$ (Proposition~\ref{prop:single-path}), which is strictly less than the consensus signal $I(D; C_K)$ for $K \ge 2$.
This negative result motivates \method{}'s fundamental departure: rather than extracting a weak signal from one path, restructure the computation to produce a multi-path signal.

% ---------------------------------------------------------------------------
\subsection{Self-Consistency and Multi-Path Reasoning}
\label{sec:related:sc}
% ---------------------------------------------------------------------------

Self-consistency~\citep{wang2023selfconsistency} improves accuracy by generating $K$ complete reasoning traces and returning the majority-vote answer.
Best-of-$N$ with reward models~\citep{cobbe2021gsm8k,lightman2023prm} samples $N$ solutions and selects the highest-scoring one according to a learned verifier.
MCTS-based reasoning methods~\citep{hao2023reasoning,xie2024selfevaluation,qi2024mutual} build search trees over reasoning steps, using rollout-based value estimates to guide expansion.
These approaches all use multiple paths to improve \emph{answer quality}---but they allocate compute \emph{uniformly} across instances, spending the same $K$ full-budget traces on every question regardless of difficulty.
\method{} draws on the same insight that multi-path agreement correlates with correctness, but uses consensus as a \emph{difficulty signal for adaptive routing}, not merely as an aggregation mechanism.
Concretely, \method{} runs $K$ \emph{short} probes (not full-budget traces), uses the agreement pattern to \emph{decide} whether additional compute is warranted, and invests that compute \emph{selectively} on hard instances.
This yields a fundamentally different cost profile: self-consistency pays $K \times B_{\mathrm{full}}$ tokens per question; \method{} pays $K \times B_{\mathrm{probe}}$ for easy questions and scales up only where consensus indicates difficulty.

% ---------------------------------------------------------------------------
\subsection{Speculative Decoding}
\label{sec:related:spec}
% ---------------------------------------------------------------------------

Speculative decoding~\citep{leviathan2023fast,chen2023accelerating} accelerates autoregressive generation by using a small, fast draft model to propose token sequences that a larger target model then verifies in parallel.
The key insight is that \emph{restructuring the computation}---from sequential single-model generation to a draft-then-verify pipeline---can reduce latency without changing the output distribution.
This paradigm has been extended to tree-structured speculation~\citep{miao2024specinfer}, self-speculation within a single model~\citep{zhang2024draft}, and cascade routing across model sizes~\citep{chen2023frugalgpt}.

\method{} draws a deliberate analogy to speculative decoding, applied at the level of \emph{reasoning} rather than token generation.
Where speculative decoding uses a cheap draft model to identify tokens the target model would likely accept, \method{} uses cheap parallel probes to identify questions the model can likely answer without extended deliberation.
Where speculative decoding falls back to the target model for tokens the draft model gets wrong, \method{} falls back to full deliberation for questions where the probes disagree.
In both cases, the approach exploits a structural asymmetry---most tokens are easy to predict, most questions are easy to solve---to reduce expected cost while preserving or improving quality, without modifying the underlying model.


% ---------------------------------------------------------------------------
\subsection{Test-Time Compute Scaling}
\label{sec:related:ttc}
% ---------------------------------------------------------------------------

The observation that ``thinking longer'' improves accuracy has motivated a new class of reasoning-specialized LLMs.
OpenAI o1~\citep{openai2024o1} first demonstrated that training models for extended chain-of-thought generation yields substantial gains on mathematical and scientific benchmarks.
QwQ~\citep{qwq2024} and DeepSeek-R1~\citep{deepseekr1} further showed that reinforcement-learning-based training for long reasoning produces strong open-weight reasoners.
\citet{snell2024scaling} formalized this as a \emph{test-time compute scaling law}: accuracy improves log-linearly with token budget on mathematical tasks.
However, this scaling is far from free---\citet{chen2024overthinking} and our own analysis (\S\ref{sec:intro}) show that accuracy \emph{degrades} on a non-trivial fraction of instances when the budget is increased, a phenomenon termed ``overthinking.''
Budget forcing~\citep{muennighoff2025s1} represents an early attempt at controlling reasoning length by truncating the thinking trace at a fixed token limit and appending a wrap-up prompt.
While effective as a diagnostic tool, budget forcing applies a uniform budget to all instances and thus cannot exploit the heterogeneity of problem difficulty---precisely the gap \method{} is designed to fill.
The ``sweet spot'' analysis of~\citet{chen2024overthinking} confirms that every instance has an optimal budget, but no fixed allocation can simultaneously serve easy and hard problems.

% ---------------------------------------------------------------------------
\subsection{Adaptive Budget Allocation}
\label{sec:related:adaptive}
% ---------------------------------------------------------------------------

A growing literature seeks to allocate reasoning compute \emph{per instance}, which we organize into three families.

\paragraph{Training-based methods.}
Several approaches modify the model or its training objective to produce difficulty-calibrated traces.
AdaThink~\citep{adathink2025} shapes rewards via GRPO to encourage the model to terminate reasoning early when confident, achieving budget savings without retraining from scratch.
AdaToken~\citep{adatoken2025} frames budget allocation as an online bandit problem and provides regret bounds, but requires online interaction during inference.
SelfBudgeter~\citep{liu2025selfbudgeter} trains the model to predict its own required budget before generation, enforcing adherence via budget-guided GRPO.
TALE~\citep{han2025tale} includes the target token budget in the prompt and trains via SFT/DPO to compress reasoning.
While these methods can be effective, they require model fine-tuning and are thus coupled to specific model weights and training data---a significant deployment constraint that \method{} avoids entirely.

\paragraph{Confidence-based methods.}
An alternative family uses the model's self-reported or estimated confidence to trigger early exit.
Satisficing Reasoning~\citep{luo2025satisficing} applies a confidence threshold to halt generation when the model deems its current answer sufficient.
Certaindex~\citep{luo2025certaindex} and Dynasor~\citep{fu2025dynasor} monitor uncertainty during generation and route to shorter or longer reasoning accordingly.
These methods are training-free but rely on a \emph{single-path} confidence signal, which we show is fundamentally limited: the model's self-assessed confidence is poorly calibrated, and single-trace uncertainty estimates achieve Spearman $|\rho| < 0.2$ with true problem difficulty (\S\ref{sec:theory}).

\paragraph{Training-free injection methods.}
RTB~\citep{hou2025thinkless} and Think Less Think Better~\citep{hou2025thinkless} inject wrap-up tokens (e.g., ``\texttt{Wait, let me reconsider\ldots}'') into the reasoning trace to steer the model toward shorter or longer deliberation without any training.
In our comprehensive evaluation (\S\ref{sec:intro}), we tested all major categories of single-path controllers---including feature-based, uncertainty-based, and injection-based methods---and found that \textbf{none outperformed the fixed-budget baseline}, with degradations of 5--7\,pp.
The root cause is informational: \emph{any} difficulty signal extracted from a single reasoning trace is bounded in mutual information with problem difficulty by $I(D; X_1)$ (Proposition~\ref{prop:single-path}), which is strictly less than the consensus signal $I(D; C_K)$ for $K \ge 2$.
This negative result motivates \method{}'s fundamental departure: rather than extracting a weak signal from one path, restructure the computation to produce a multi-path signal.

% ---------------------------------------------------------------------------
\subsection{Self-Consistency and Multi-Path Reasoning}
\label{sec:related:sc}
% ---------------------------------------------------------------------------

Self-consistency~\citep{wang2023selfconsistency} improves accuracy by generating $K$ complete reasoning traces and returning the majority-vote answer.
Best-of-$N$ with reward models~\citep{cobbe2021gsm8k,lightman2023prm} samples $N$ solutions and selects the highest-scoring one according to a learned verifier.
MCTS-based reasoning methods~\citep{hao2023reasoning,xie2024selfevaluation,qi2024mutual} build search trees over reasoning steps, using rollout-based value estimates to guide expansion.
These approaches all use multiple paths to improve \emph{answer quality}---but they allocate compute \emph{uniformly} across instances, spending the same $K$ full-budget traces on every question regardless of difficulty.
\method{} draws on the same insight that multi-path agreement correlates with correctness, but uses consensus as a \emph{difficulty signal for adaptive routing}, not merely as an aggregation mechanism.
Concretely, \method{} runs $K$ \emph{short} probes (not full-budget traces), uses the agreement pattern to \emph{decide} whether additional compute is warranted, and invests that compute \emph{selectively} on hard instances.
This yields a fundamentally different cost profile: self-consistency pays $K \times B_{\mathrm{full}}$ tokens per question; \method{} pays $K \times B_{\mathrm{probe}}$ for easy questions and scales up only where consensus indicates difficulty.

% ---------------------------------------------------------------------------
\subsection{Speculative Decoding}
\label{sec:related:spec}
% ---------------------------------------------------------------------------

Speculative decoding~\citep{leviathan2023fast,chen2023accelerating} accelerates autoregressive generation by using a small, fast draft model to propose token sequences that a larger target model then verifies in parallel.
The key insight is that \emph{restructuring the computation}---from sequential single-model generation to a draft-then-verify pipeline---can reduce latency without changing the output distribution.
This paradigm has been extended to tree-structured speculation~\citep{miao2024specinfer}, self-speculation within a single model~\citep{zhang2024draft}, and cascade routing across model sizes~\citep{chen2023frugalgpt}.

\method{} draws a deliberate analogy to speculative decoding, applied at the level of \emph{reasoning} rather than token generation.
Where speculative decoding uses a cheap draft model to identify tokens the target model would likely accept, \method{} uses cheap parallel probes to identify questions the model can likely answer without extended deliberation.
Where speculative decoding falls back to the target model for tokens the draft model gets wrong, \method{} falls back to full deliberation for questions where the probes disagree.
In both cases, the approach exploits a structural asymmetry---most tokens are easy to predict, most questions are easy to solve---to reduce expected cost while preserving or improving quality, without modifying the underlying model.


% ---------------------------------------------------------------------------
\subsection{Test-Time Compute Scaling}
\label{sec:related:ttc}
% ---------------------------------------------------------------------------

The observation that ``thinking longer'' improves accuracy has motivated a new class of reasoning-specialized LLMs.
OpenAI o1~\citep{openai2024o1} first demonstrated that training models for extended chain-of-thought generation yields substantial gains on mathematical and scientific benchmarks.
QwQ~\citep{qwq2024} and DeepSeek-R1~\citep{deepseekr1} further showed that reinforcement-learning-based training for long reasoning produces strong open-weight reasoners.
\citet{snell2024scaling} formalized this as a \emph{test-time compute scaling law}: accuracy improves log-linearly with token budget on mathematical tasks.
However, this scaling is far from free---\citet{chen2024overthinking} and our own analysis (\S\ref{sec:intro}) show that accuracy \emph{degrades} on a non-trivial fraction of instances when the budget is increased, a phenomenon termed ``overthinking.''
Budget forcing~\citep{muennighoff2025s1} represents an early attempt at controlling reasoning length by truncating the thinking trace at a fixed token limit and appending a wrap-up prompt.
While effective as a diagnostic tool, budget forcing applies a uniform budget to all instances and thus cannot exploit the heterogeneity of problem difficulty---precisely the gap \method{} is designed to fill.
The ``sweet spot'' analysis of~\citet{chen2024overthinking} confirms that every instance has an optimal budget, but no fixed allocation can simultaneously serve easy and hard problems.

% ---------------------------------------------------------------------------
\subsection{Adaptive Budget Allocation}
\label{sec:related:adaptive}
% ---------------------------------------------------------------------------

A growing literature seeks to allocate reasoning compute \emph{per instance}, which we organize into three families.

\paragraph{Training-based methods.}
Several approaches modify the model or its training objective to produce difficulty-calibrated traces.
AdaThink~\citep{adathink2025} shapes rewards via GRPO to encourage the model to terminate reasoning early when confident, achieving budget savings without retraining from scratch.
AdaToken~\citep{adatoken2025} frames budget allocation as an online bandit problem and provides regret bounds, but requires online interaction during inference.
SelfBudgeter~\citep{liu2025selfbudgeter} trains the model to predict its own required budget before generation, enforcing adherence via budget-guided GRPO.
TALE~\citep{han2025tale} includes the target token budget in the prompt and trains via SFT/DPO to compress reasoning.
While these methods can be effective, they require model fine-tuning and are thus coupled to specific model weights and training data---a significant deployment constraint that \method{} avoids entirely.

\paragraph{Confidence-based methods.}
An alternative family uses the model's self-reported or estimated confidence to trigger early exit.
Satisficing Reasoning~\citep{luo2025satisficing} applies a confidence threshold to halt generation when the model deems its current answer sufficient.
Certaindex~\citep{luo2025certaindex} and Dynasor~\citep{fu2025dynasor} monitor uncertainty during generation and route to shorter or longer reasoning accordingly.
These methods are training-free but rely on a \emph{single-path} confidence signal, which we show is fundamentally limited: the model's self-assessed confidence is poorly calibrated, and single-trace uncertainty estimates achieve Spearman $|\rho| < 0.2$ with true problem difficulty (\S\ref{sec:theory}).

\paragraph{Training-free injection methods.}
RTB~\citep{hou2025thinkless} and Think Less Think Better~\citep{hou2025thinkless} inject wrap-up tokens (e.g., ``\texttt{Wait, let me reconsider\ldots}'') into the reasoning trace to steer the model toward shorter or longer deliberation without any training.
In our comprehensive evaluation (\S\ref{sec:intro}), we tested all major categories of single-path controllers---including feature-based, uncertainty-based, and injection-based methods---and found that \textbf{none outperformed the fixed-budget baseline}, with degradations of 5--7\,pp.
The root cause is informational: \emph{any} difficulty signal extracted from a single reasoning trace is bounded in mutual information with problem difficulty by $I(D; X_1)$ (Proposition~\ref{prop:single-path}), which is strictly less than the consensus signal $I(D; C_K)$ for $K \ge 2$.
This negative result motivates \method{}'s fundamental departure: rather than extracting a weak signal from one path, restructure the computation to produce a multi-path signal.

% ---------------------------------------------------------------------------
\subsection{Self-Consistency and Multi-Path Reasoning}
\label{sec:related:sc}
% ---------------------------------------------------------------------------

Self-consistency~\citep{wang2023selfconsistency} improves accuracy by generating $K$ complete reasoning traces and returning the majority-vote answer.
Best-of-$N$ with reward models~\citep{cobbe2021gsm8k,lightman2023prm} samples $N$ solutions and selects the highest-scoring one according to a learned verifier.
MCTS-based reasoning methods~\citep{hao2023reasoning,xie2024selfevaluation,qi2024mutual} build search trees over reasoning steps, using rollout-based value estimates to guide expansion.
These approaches all use multiple paths to improve \emph{answer quality}---but they allocate compute \emph{uniformly} across instances, spending the same $K$ full-budget traces on every question regardless of difficulty.
\method{} draws on the same insight that multi-path agreement correlates with correctness, but uses consensus as a \emph{difficulty signal for adaptive routing}, not merely as an aggregation mechanism.
Concretely, \method{} runs $K$ \emph{short} probes (not full-budget traces), uses the agreement pattern to \emph{decide} whether additional compute is warranted, and invests that compute \emph{selectively} on hard instances.
This yields a fundamentally different cost profile: self-consistency pays $K \times B_{\mathrm{full}}$ tokens per question; \method{} pays $K \times B_{\mathrm{probe}}$ for easy questions and scales up only where consensus indicates difficulty.

% ---------------------------------------------------------------------------
\subsection{Speculative Decoding}
\label{sec:related:spec}
% ---------------------------------------------------------------------------

Speculative decoding~\citep{leviathan2023fast,chen2023accelerating} accelerates autoregressive generation by using a small, fast draft model to propose token sequences that a larger target model then verifies in parallel.
The key insight is that \emph{restructuring the computation}---from sequential single-model generation to a draft-then-verify pipeline---can reduce latency without changing the output distribution.
This paradigm has been extended to tree-structured speculation~\citep{miao2024specinfer}, self-speculation within a single model~\citep{zhang2024draft}, and cascade routing across model sizes~\citep{chen2023frugalgpt}.

\method{} draws a deliberate analogy to speculative decoding, applied at the level of \emph{reasoning} rather than token generation.
Where speculative decoding uses a cheap draft model to identify tokens the target model would likely accept, \method{} uses cheap parallel probes to identify questions the model can likely answer without extended deliberation.
Where speculative decoding falls back to the target model for tokens the draft model gets wrong, \method{} falls back to full deliberation for questions where the probes disagree.
In both cases, the approach exploits a structural asymmetry---most tokens are easy to predict, most questions are easy to solve---to reduce expected cost while preserving or improving quality, without modifying the underlying model.
